% !TEX root = ../main.tex
\begin{ouverture}[newton]{15}{Travail et énergie cinétique}{Partie 3 · L'énergie : conversions et transferts}
{Énergie cinétique\enspace·\enspace Travail d'une force\enspace·\enspace Travail du poids\enspace·\enspace Frottements\enspace·\enspace Théorème de l'énergie cinétique}
\objectif{Utiliser l'expression de l'énergie cinétique d'un système modélisé par un point matériel.}
\objectif{Utiliser l'expression du travail $W_{\mathrm{AB}}(\vec{F})=\vec{F}\cdot\overrightarrow{\mathrm{AB}}$ d'une force constante.}
\objectif{Calculer le travail du poids et celui d'une force de frottement d'intensité constante sur une trajectoire rectiligne.}
\objectif{Énoncer et exploiter le théorème de l'énergie cinétique.}
\prerequis{Énergie cinétique, conversions d'énergie (collège) ; forces usuelles, poids $P=m\,g$.}
\prerequis{Vecteurs, angles ; cosinus d'un angle (mathématiques) ; vitesse en \unit{m.s^{-1}}.}
\end{ouverture}
\essentiel
\essSection{1}{Énergie cinétique}
\begin{definition}[Énergie cinétique]
L'énergie cinétique d'un système de masse $m$ (en kg), modélisé par un point de vitesse $v$ (en \unit{m.s^{-1}}), vaut
\[E_\mathrm{c}=\frac{1}{2}\,m\,v^2\qquad\text{(en joules, J).}\]
Elle est proportionnelle au carré de la vitesse : à vitesse double, énergie quadruple. Une voiture de \qty{1000}{kg} à \qty{72}{km/h} (\qty{20}{m.s^{-1}})
possède \qty{2,0e5}{J}.
\end{definition}
\essSection{2}{Travail d'une force constante}
\begin{definition}[Travail d'une force constante]
Le travail d'une force constante $\vec{F}$ dont le point d'application se déplace de A à B est le produit scalaire
\[W_{\mathrm{AB}}(\vec{F})=\vec{F}\cdot\overrightarrow{\mathrm{AB}}=F\times\mathrm{AB}\times\cos\alpha\qquad\text{(en J ; $F$ en N, AB en m),}\]
où $\alpha$ est l'angle entre $\vec{F}$ et $\overrightarrow{\mathrm{AB}}$. Le travail est \textbf{moteur} ($W>0$) si $\alpha<\ang{90}$, \textbf{nul} si la force est perpendiculaire
au déplacement, \textbf{résistant} ($W<0$) si $\alpha>\ang{90}$. Exemple : \qty{100}{N} à \ang{60} sur \qty{10}{m} donnent \qty{500}{J}.
\end{definition}
\schema{15-travail}
\begin{propriete}[Deux forces usuelles]
\textbf{Poids} : $W_{\mathrm{AB}}(\vec{P})=m\,g\,(z_\mathrm{A}-z_\mathrm{B})$, avec un axe $z$ vertical vers le haut. Il ne dépend que des altitudes de départ et
d'arrivée, pas du chemin suivi : on dit que le poids est une force \textbf{conservative}.\par
\textbf{Frottement} d'intensité constante $f$, opposé au mouvement, sur une trajectoire rectiligne : $W_{\mathrm{AB}}(\vec{f})=-f\times\mathrm{AB}$. Toujours
résistant, il dépend de la longueur du trajet : force \textbf{non conservative}.
\end{propriete}
\par\penalty0
\schema{15-poids}
\essSection{3}{Théorème de l'énergie cinétique}
\begin{loi}[Théorème de l'énergie cinétique]
Entre deux positions A et B, la variation de l'énergie cinétique d'un système est égale à la somme des travaux des forces
qui s'exercent sur lui :
\[E_\mathrm{c}(\mathrm{B})-E_\mathrm{c}(\mathrm{A})=\sum W_{\mathrm{AB}}(\vec{F}).\]
Les forces motrices augmentent l'énergie cinétique, les forces résistantes la diminuent. Exemple : un objet lâché sans vitesse
qui tombe de \qty{5,0}{m} sans frottement atteint $v=\sqrt{2\,g\,h}=\qty{9,9}{m.s^{-1}}$, quelle que soit sa masse.
\end{loi}
\par\penalty0
\schema{15-tec}
\begin{python}
m, g, v = 0.50, 9.81, 0.0          # kg, N/kg, m/s : objet lache
for h in [1, 2, 3, 4, 5]:          # hauteur de chute en m
    v = (2 * g * h) ** 0.5         # Ec = m g h (theoreme)
    print(h, 'm :', round(v, 1), 'm/s ; Ec =', round(0.5 * m * v**2, 1), 'J')
\end{python}
\begin{console}
1 m : 4.4 m/s ; Ec = 4.9 J
2 m : 6.3 m/s ; Ec = 9.8 J
3 m : 7.7 m/s ; Ec = 14.7 J
4 m : 8.9 m/s ; Ec = 19.6 J
5 m : 9.9 m/s ; Ec = 24.5 J
\end{console}
\noindent\capacite{Capacité numérique} L'énergie cinétique gagnée est proportionnelle à la hauteur de chute, mais la vitesse ne l'est pas : il faut
quatre fois plus de hauteur pour doubler la vitesse.
\begin{attention}
Le travail d'une force n'est pas une force : c'est une énergie, en joules. Ne pas oublier les forces dont le travail est nul (réaction
d'un support sans frottement) dans l'inventaire, et utiliser la vitesse en \unit{m.s^{-1}} dans $E_\mathrm{c}$.
\end{attention}
\begin{savoirfaire}[Appliquer le théorème de l'énergie cinétique]
\begin{enumerate}
\item Définir le système et les positions initiale A et finale B ; faire l'inventaire des forces et un schéma.
\item Calculer le travail de chaque force entre A et B, avec son signe.
\item Écrire $E_\mathrm{c}(\mathrm{B})-E_\mathrm{c}(\mathrm{A})=\sum W_{\mathrm{AB}}$ et isoler l'inconnue : vitesse, distance ou force.
\end{enumerate}
\end{savoirfaire}
\begin{experience}[Étudier une chute ou une glissade]
On filme la descente d'un mobile sur un rail incliné. Un logiciel de pointage donne la vitesse en différents points ; on compare
l'énergie cinétique gagnée au travail du poids et l'on en déduit le travail des frottements.
\end{experience}
\begin{rediger}
Préciser le système, le référentiel et les deux positions ; écrire chaque travail littéralement avec son signe, puis le théorème,
avant de passer aux valeurs numériques.
\end{rediger}
\exercices

\rubrique{Découverte}
\begin{blocLivre}{ex-15-1}
\begin{exercice}[1]{Calculer des énergies cinétiques}
\begin{questions}
\item Calculer l'énergie cinétique d'une voiture de \qty{1200}{kg} roulant à \qty{90}{km/h}, puis à \qty{130}{km/h}.
\item Par quel facteur l'énergie cinétique est-elle multipliée ? Le comparer au rapport des vitesses.
\item Calculer l'énergie cinétique d'une balle de tennis de \qty{58}{g} à \qty{50}{m.s^{-1}} et celle d'un piéton de \qty{70}{kg} marchant à \qty{1,4}{m.s^{-1}}.
\item À quelle vitesse la voiture aurait-elle la même énergie cinétique que la balle ?
\end{questions}
\end{exercice}
\end{blocLivre}
\begin{blocLivre}{ex-15-2}
\begin{exercice}[1]{Tirer une luge}
\schema{15-ex02}
\begin{questions}
\item Calculer le travail de la force exercée par l'enfant sur \qty{20}{m}.
\item Que vaudrait ce travail si la corde était horizontale ?
\item Représenter la composante de $\vec{F}$ qui travaille et calculer sa valeur.
\item Quel est le travail du poids de la luge ? Justifier.
\end{questions}
\end{exercice}
\end{blocLivre}
\begin{blocLivre}{ex-15-3}
\begin{exercice}[1]{Le travail du poids}
\begin{donnees}
$g=\qty{9,81}{N.kg^{-1}}$.
\end{donnees}
\begin{questions}
\item Une balle de \qty{0,50}{kg} tombe de \qty{2,0}{m}. Calculer le travail de son poids. Est-il moteur ou résistant ?
\item Une personne de \qty{60}{kg} monte un étage de \qty{3,0}{m}. Calculer le travail de son poids.
\item Un randonneur de \qty{70}{kg} gravit un sommet situé \qty{800}{m} plus haut que son point de départ. Calculer le travail de son poids.
\item Le travail est-il le même s'il emprunte un sentier en lacets plutôt qu'un chemin direct ? Justifier.
\end{questions}
\end{exercice}
\end{blocLivre}
\begin{blocLivre}{ex-15-4}
\begin{exercice}[1]{Moteur, résistant ou nul ?}
\schema{15-ex04}
\begin{questions}
\item Pour chaque situation, indiquer l'angle entre la force et le déplacement.
\item En déduire si le travail de la force est moteur, résistant ou nul.
\item Dans quelle situation la force fait-elle augmenter l'énergie cinétique ?
\item Pourquoi le poids d'un objet qui glisse sur une table horizontale ne travaille-t-il pas ?
\end{questions}
\end{exercice}
\end{blocLivre}
\begin{blocLivre}{ex-15-5}
\begin{exercice}[1]{Une chute libre}
On lâche une bille de \qty{0,20}{kg} sans vitesse initiale d'une hauteur de \qty{5,0}{m}. On néglige l'action de l'air.
\begin{donnees}
$g=\qty{9,81}{N.kg^{-1}}$.
\end{donnees}
\begin{questions}
\item Faire l'inventaire des forces et calculer le travail du poids pendant la chute.
\item Appliquer le théorème de l'énergie cinétique et calculer la vitesse de la bille à l'arrivée au sol.
\item La vitesse dépend-elle de la masse de la bille ?
\item Quelle serait la vitesse après une chute de \qty{20}{m} ? Commenter.
\end{questions}
\end{exercice}
\end{blocLivre}
\rubrique{Entraînement}
\begin{blocLivre}{ex-15-6}
\begin{exercice}[2]{Le travail selon l'angle, en Python}
\begin{python}
from math import cos, radians
F, d = 100, 10                     # N, m
for angle in [0, 30, 60, 90, 120, 180]:
    W = F * d * cos(radians(angle))
    print(angle, 'degres :', round(W), 'J')
\end{python}
\begin{console}
0 degres : 1000 J
30 degres : 866 J
60 degres : 500 J
90 degres : 0 J
120 degres : -500 J
180 degres : -1000 J
\end{console}
\begin{questions}
\item \capacite{Capacité numérique}Quelle relation le programme utilise-t-il ? Pourquoi faut-il la fonction \code{radians} ?
\item Pour quels angles le travail est-il moteur ? Résistant ? Nul ?
\item Vérifier le résultat affiché pour \ang{60}.
\item Donner un exemple concret de force faisant un angle de \ang{180} avec le déplacement.
\end{questions}
\end{exercice}
\end{blocLivre}
\begin{blocLivre}{ex-15-7}
\begin{exercice}[2]{Le freinage d'une voiture}
Une voiture de \qty{1000}{kg} roulant à \qty{90}{km/h} freine en ligne droite et s'arrête sur \qty{50}{m}. On suppose la force de freinage constante.
\begin{questions}
\item Calculer la variation de l'énergie cinétique de la voiture.
\item Quelles forces s'exercent sur la voiture ? Lesquelles travaillent ?
\item Appliquer le théorème de l'énergie cinétique pour calculer la force de freinage.
\item Avec la même force, sur quelle distance la voiture s'arrêterait-elle depuis \qty{130}{km/h} ?
\end{questions}
\end{exercice}
\end{blocLivre}
\begin{blocLivre}{ex-15-8}
\begin{exercice}[2]{La descente du skieur}
\schema{15-ex08}
\begin{donnees}
$g=\qty{9,81}{N.kg^{-1}}$.
\end{donnees}
\begin{questions}
\item Faire l'inventaire des forces qui s'exercent sur le skieur et les représenter.
\item Calculer le travail du poids entre A et B. En déduire la vitesse en B en l'absence de frottement.
\item La vitesse mesurée en B vaut \qty{18}{m.s^{-1}}. Calculer le travail des frottements.
\item En déduire la valeur de la force de frottement, supposée constante.
\end{questions}
\end{exercice}
\end{blocLivre}
\begin{blocLivre}{ex-15-9}
\begin{exercice}[2]{Lancer dans toutes les directions}
Une balle est lancée à \qty{10}{m.s^{-1}} depuis une hauteur de \qty{1,5}{m}, horizontalement, vers le haut ou vers le bas. On néglige l'air.
\begin{donnees}
$g=\qty{9,81}{N.kg^{-1}}$.
\end{donnees}
\begin{questions}
\item Appliquer le théorème de l'énergie cinétique entre le lancer et l'arrivée au sol.
\item Montrer que $v_\mathrm{sol}=\sqrt{v_0^2+2\,g\,h}$ et calculer cette vitesse.
\item La direction du lancer a-t-elle une influence sur la vitesse d'arrivée au sol ?
\item Calculer la vitesse au sol si la balle est lancée depuis \qty{2,0}{m}.
\end{questions}
\end{exercice}
\end{blocLivre}
\begin{blocLivre}{ex-15-10}
\begin{exercice}[2]{Le chariot de supermarché}
\schema{15-ex10}
On pousse un chariot de \qty{30}{kg}, initialement immobile, avec une force de \qty{40}{N}.
\begin{questions}
\item Calculer le travail de la force de poussée sur \qty{15}{m}.
\item Calculer le travail des frottements. Quels sont les travaux du poids et de la réaction du sol ?
\item Appliquer le théorème de l'énergie cinétique et calculer la vitesse finale du chariot.
\item Pourquoi est-il plus efficace de pousser horizontalement ?
\end{questions}
\end{exercice}
\end{blocLivre}
\begin{blocLivre}{ex-15-11}
\begin{exercice}[2]{Descendre un col à vitesse constante}
Une voiture de \qty{1500}{kg} descend un col à vitesse constante : sur \qty{200}{m} de route, elle perd \qty{16}{m} d'altitude.
\begin{donnees}
$g=\qty{9,81}{N.kg^{-1}}$.
\end{donnees}
\begin{questions}
\item Que vaut la variation d'énergie cinétique de la voiture ?
\item Calculer le travail du poids sur ce trajet.
\item En déduire le travail des forces de freinage (freins et frottements), puis leur valeur moyenne.
\item Que devient l'énergie correspondante ? Pourquoi conseille-t-on d'utiliser le frein moteur en descente ?
\end{questions}
\end{exercice}
\end{blocLivre}
\rubrique{Objectif DS}
\begin{blocLivre}{ex-15-12}
\begin{exercice}[3]{Le tir à l'arc}
\begin{docu}{1}{La corde de l'arc}
L'archer tend la corde en reculant la flèche de \qty{50}{cm}. Lorsqu'il la lâche, la corde exerce sur la flèche une force moyenne de \qty{150}{N},
dans le sens du déplacement, sur ces \qty{50}{cm}.
\end{docu}
\begin{docu}{2}{La flèche}
Masse : \qty{20}{g}. Le travail du poids pendant la détente est négligeable.
\end{docu}
\begin{donnees}
$g=\qty{9,81}{N.kg^{-1}}$.
\end{donnees}
\begin{questions}
\item Calculer le travail de la force exercée par la corde.
\item Appliquer le théorème de l'énergie cinétique et calculer la vitesse de la flèche, en \unit{m.s^{-1}} puis en \unit{km/h}.
\item \initiative Sans frottement, quelle altitude atteindrait une flèche tirée verticalement ?
\item En réalité, la flèche monte beaucoup moins haut. Pourquoi ?
\end{questions}
\end{exercice}
\end{blocLivre}
\begin{blocLivre}{ex-15-13}
\begin{exercice}[3]{La distance de freinage}
\begin{docu}{1}{Freinage d'urgence}
Lors d'un freinage d'urgence, la route exerce sur une voiture de masse $m$ une force de freinage d'intensité $f=0{,}8\,m\,g$ sur route sèche et
$f=0{,}5\,m\,g$ sur route mouillée, opposée au mouvement.
\end{docu}
\begin{donnees}
$g=\qty{9,81}{N.kg^{-1}}$.
\end{donnees}
\begin{questions}
\item Appliquer le théorème de l'énergie cinétique pour montrer que la distance de freinage vaut $d=\dfrac{v^2}{2\times0{,}8\,g}$ sur route sèche.
\item Calculer la distance de freinage à 50, 90 et \qty{130}{km/h}.
\item Que devient la distance de freinage quand la vitesse double ? La masse de la voiture intervient-elle ?
\item \initiative Calculer la distance de freinage à \qty{90}{km/h} sur route mouillée. Conclure.
\end{questions}
\end{exercice}
\end{blocLivre}
\begin{blocLivre}{ex-15-14}
\begin{exercice}[3]{Le saut à la perche}
\begin{docu}{1}{Le principe}
Le perchiste court, plante sa perche, qui se courbe puis se redresse : l'énergie cinétique de la course sert à élever le sauteur.
Au mieux, on peut considérer qu'elle est entièrement utilisée pour monter, la vitesse au sommet étant presque nulle.
\end{docu}
\begin{docu}{2}{Le sauteur}
Masse : \qty{75}{kg} ; vitesse de course juste avant le saut : \qty{9,5}{m.s^{-1}} ; son centre de gravité est alors à \qty{1,1}{m} du sol.
\end{docu}
\begin{donnees}
$g=\qty{9,81}{N.kg^{-1}}$.
\end{donnees}
\begin{questions}
\item Calculer l'énergie cinétique du perchiste avant le saut.
\item En appliquant le théorème de l'énergie cinétique entre le début du saut et le sommet, calculer l'élévation maximale de son centre de gravité.
\item En déduire la hauteur maximale atteinte par son centre de gravité.
\item \initiative Les meilleurs perchistes franchissent plus de \qty{6,2}{m}. Quelle vitesse de course faudrait-il, d'après ce modèle ? Commenter.
\end{questions}
\end{exercice}
\end{blocLivre}
\begin{blocLivre}{ex-15-15}
\begin{exercice}[3]{L'atterrissage d'un avion de ligne}
\begin{docu}{1}{Les données}
Un avion de ligne de \qty{64}{t} se pose à \qty{70}{m.s^{-1}} et s'arrête après \qty{1500}{m} de roulage sur une piste horizontale. Pour freiner, il utilise ses
freins et ses inverseurs de poussée, qui dirigent vers l'avant le jet des moteurs.
\end{docu}
\begin{donnees}
$g=\qty{9,81}{N.kg^{-1}}$.
\end{donnees}
\begin{questions}
\item Calculer l'énergie cinétique de l'avion au moment de l'atterrissage.
\item En appliquant le théorème de l'énergie cinétique, calculer la force de freinage moyenne. La comparer au poids de l'avion.
\item Quels sont les travaux du poids et de la réaction de la piste ?
\item \initiative Les inverseurs de poussée assurent \qty{40}{\percent} de la force de freinage. Sur quelle distance l'avion s'arrêterait-il sans eux, avec les
seuls freins ?
\end{questions}
\end{exercice}
\end{blocLivre}
